\subsection{Arithmetic equivalence from sparse Li coefficients}

Let $d_n=\lambda_n(K)-\lambda_n(L)$ for two number fields $K,L$.
The following inverse result requires no critical-line hypothesis.

\paragraph{Theorem.}
If $K,L$ are arithmetically inequivalent, meaning that
$\zeta_K\ne\zeta_L$, then
\[
 \lim_{H\to\infty}\ \inf_{A\ge1}
   \frac1H\sum_{n=A}^{A+H-1}|d_n|^2=+\infty,
 \qquad A,H\in\mathbb N .
\]
Thus even a bounded mean square on one sequence of windows of increasing
length forces $\zeta_K=\zeta_L$. The windows may
be arbitrarily far apart. For example the windows starting at
$A_j=2^{2^{j^2}}$ and of length $j$ use only $O(1+\log\log N)$ indices
up to $N$. Equality of zeta functions gives $d_n=0$ for every $n$.

\paragraph{Proof.}
Write $m_w$ for the signed difference of the multiplicities of the
transformed zero $w=1-1/\rho$. Taking second differences of the two
symmetric zero sums gives
\[
 a_n:=-(d_{n+2}-2d_{n+1}+d_n)=\sum_w c_w w^n,\qquad
 c_w=m_w(1-w)^2,\qquad \sum_w|c_w|<\infty.
\]
Indeed, $|1-w|=|\rho|^{-1}$ and the zero count gives convergence
of the square sum. All $|w|$ are bounded. Suppose, towards the
contrapositive, that there are windows of length $H\to\infty$
and start $A$ on which the mean square of $d_n$ is at most $C^2$.
Then $|d_n|\le C\sqrt H$ throughout each window.

If an unmatched $w$ has modulus above $1$, let $R$ be the largest
such modulus. It is attained, since $w\to1$ at increasing height.
There are finitely many distinct unmatched points $w_1,\ldots,w_t$
with this modulus, and the other unmatched moduli are bounded by
some $r<R$. Put $z_i=w_i/R$. For $N=A+H-t-2$ and $0\le h<t$,
\[
 R^{-N-h}a_{N+h}
 =\sum_{i=1}^t z_i^h c_{w_i}z_i^N +O((r/R)^{N+h}).
\]
The left sides tend to zero since $N\ge H-t-1$ and
$|a_{N+h}|\le4C\sqrt H$. The fixed Vandermonde matrix
$(z_i^h)_{h,i}$ is invertible. Consequently $c_{w_i}z_i^N\to0$,
contradicting $|c_{w_i}|>0$. Thus all outer points match.
Reflection $\rho\mapsto1-\rho$ induces $w\mapsto1/w$ and preserves
multiplicity, so all inner points match as well.

All remaining terms have $|w|=1$. Fix such a $z$.
The average of $a_n z^{-n}$ over $A\le n\le A+H-3$, normalized
by $H$, tends to $c_z$. Indeed, each other mass has geometric
average bounded in modulus by $\min(1,2/(H|1-w/z|))$, uniformly
in $A$, and the masses are absolutely summable. Shifting the two
indices in the second difference now gives
\[
 \frac1H\sum_{n=A}^{A+H-3}a_nz^{-n}
 =-(z-1)^2\frac1H\sum_{n=A}^{A+H-1}d_nz^{-n}
       +O(C/\sqrt H).
\]
Every omitted endpoint is bounded by $C\sqrt H$. Since $z\ne1$,
the normalized Fourier averages of $d_n$ tend to $-m_z$.
For any finite set of distinct unit points, their exponentials on
the window have normalized Gram matrix tending to the identity, by
the same geometric-sum estimate. The projection inequality
(or Cauchy--Schwarz applied to their linear combinations) therefore
gives $\sum_z |m_z|^2\le C^2$ on every such finite set.
The multiplicities are integers. Only finitely many can differ.

The pole-cleared completions $\xi_K,\xi_L$ are entire of order one.
Hadamard factorization gives $\xi_K/\xi_L=e^{a+bs}P(s)/Q(s)$
for polynomials $P,Q$. Set $Z(s)=\zeta_K(s)/\zeta_L(s)$.
The gamma factors of a number field are products of
$\Gamma_{\mathbb R}(s)$ and $\Gamma_{\mathbb R}(s+1)$ with
$\Gamma_{\mathbb R}(s)=\pi^{-s/2}\Gamma(s/2)$. Their two-step
quotients are rational functions, since
$\Gamma_{\mathbb R}(s+2)/\Gamma_{\mathbb R}(s)=s/(2\pi)$.
The discriminant exponential contributes a constant under the same
shift. Hence $Z(s+2)/Z(s)$ is rational. On the positive real axis,
$\zeta_K(s)=1+O_K(2^{-s})$ by absolute convergence of the
Dirichlet series at $2$, and likewise for $L$. Thus
$Z(s+2)/Z(s)=1+O_{K,L}(2^{-s})$. A rational function with this
decay towards $1$ must be identically $1$. For real $s>1$,
$Z(s)=Z(s+2m)\to1$, proving equality of the two zeta functions.

This identifies a zeta function up to equality and allows nonisomorphic
arithmetically equivalent fields. The theorem uses the classical Li
zero sum and a uniform Fourier average on long integer windows.
It supplies an infinite-data test and gives no effective length at
which a specified pair of fields can be distinguished.
